% ================================================================
% ==== FILE: content/k_levels/klevels_master.tex
% ================================================================



\paragraph{Canonical tuple projection note.}
Local K-level displays in this module are projections of the canonical public tuple \(K=(\Omega,\partial\Omega,A,\Theta,P,J,C,k,M)\). When a display omits \(M\), reorders components, or uses level-indexed shorthand, the omitted embedding context is inherited from the surrounding level discussion rather than introduced as a competing definition.

\section{Erweiterte Module und stufenübergreifende Strukturen}
\label{sec:klevels-master}

Die Ontologie der Kontinua ist als hierarchische Abfolge von organisiert
Ebenen
\[
K_0, K_1, \dots, K_{12},
\]
jedes stellt eine qualitativ unterschiedliche Strukturform dar,
Dynamik und zulässige Zustände.
Jede Ebene $K_x$ besitzt:
\begin{itemize}
    \item ein Zustandsraum $\Omega(K_x)$,
    \item eine Menge von Achsen $A(K_x)$,
    \item Potentiale $P(K_x)$,
    \item Schwellenwerte $\Theta(K_x)$,
    \item fließt $J(K_x)$,
    \item Zyklen $C(K_x)$ mit charakteristischen Zeiten $\tau(K_x)$,
    \item ein Kontinuumsmaß $k(K_x)$,
    \item und eine Grenze $\partial\Omega(K_x)$, die die Bedingungen beschreibt
          des Zusammenbruchs oder Übergangs.
\end{itemize}

Jedes $K_x$ ist in einen entsprechenden Metaraum $M_x$ eingebettet
ermittelt die zulässigen Achsen und gibt an:
\[
A(K_x) \subseteq A(M_x), \qquad
\Omega(K_x) \subseteq \Omega(M_x).
\]
Der Metaraum schränkt die Freiheitsgrade ein, die dem zur Verfügung stehen
Kontinuum und prägt dadurch seine Entwicklung.


% ================================================================
\subsection{Dimensionales Wachstum und Übergänge $K_x \to K_{x+1}$}
% ================================================================

Ein Übergang vom Niveau $K_x$ zu $K_{x+1}$ findet statt, wenn:
\begin{enumerate}
    \item Es entsteht eine neue Klasse von Unterschieden, die nicht dargestellt werden können
          innerhalb der vorhandenen Achsen $A(K_x)$;
    \item Strukturspannung überschreitet einen Schwellenwert:
    \[
        T(K_x) > \Theta_{\mathrm{dim}}(K_x);
    \]
    \item der Metaraum erfordert eine höherdimensionale Darstellung;
    \item wird der bestehende Konfigurationsraum $\Omega(K_x)$
          unzureichend oder kollabiert auf seiner Grenze $\partial\Omega(K_x)$.
\end{enumerate}

Dieser Vorgang ist irreversibel:
\[
\dim K_{x+1} \ge \dim K_x,
\]
und jede „Reduktion“ der Dimensionalität entspricht dem Tod des
Kontinuum:
\[
\Omega(K_x) = \varnothing.
\]


% ================================================================
\subsection{Universelle Struktur über alle Ebenen}
% ================================================================

Trotz der qualitativen Vielfalt der Ebenen hat jedes $K_x$ etwas gemeinsam
formale Struktur.
Die Entwicklung eines Kontinuums wird bestimmt durch:
\[
K_x(t + dt) = E\bigl(K_x(t), M_x(t)\bigr),
\]
wobei $E$ ein durch den Metaraum eingeschränkter Evolutionsoperator ist
$M_x$.
Ein Kontinuum entwickelt sich nur innerhalb seines zulässigen Bereichs
Metaraum:
\[
K_x(t+dt) \in \Omega(M_x)
\quad\text{oder die Evolution ist verboten}.
\]

Das Kontinuumsmaß ist:
\[
k(K_x) =
k_{\Omega} \cdot k_{A} \cdot k_{C} \cdot k_{J} \cdot k_{T},
\]
wo:
\begin{itemize}
    \item $k_{\Omega}$ misst das Volumen der erlaubten Zustände,
    \item $k_{A}$ zählt realisierte Achsen relativ zum Metaraum,
    \item $k_{C}$ spiegelt die Stabilität von Zyklen wider,
    \item $k_{J}$ beschreibt die Kohärenz von Flüssen,
    \item $k_{T}$ misst die strukturelle Spannung relativ zu Schwellenwerten.
\end{itemize}


% ================================================================
\subsection{Hierarchie der Ebenen $K_0 \to K_{12}$}
% ================================================================

Jede Ebene stellt eine bestimmte Organisationsschicht dar:

\begin{itemize}
    \item $K_0$ – Metadruck, Möglichkeitsbedingungen und prästrukturelle Logik;
    \item $K_1$ — kontinuierliche eindimensionale Kontinua;
    \item $K_2$ – zweidimensionale Kontinua mit internen Flüssen und Clustern;
    \item $K_3$ – chemische Konfigurationsräume und Reaktionsmannigfaltigkeiten;
    \item $K_4$ – Kontinua auf Protozellenebene mit Membranen und Grenzen;
    \item $K_5$ – neuronale Kontinua mit Erregungszyklen und Attraktoren;
    \item $K_6$ – kognitive Kontinua mit Vorhersagemodellen;
    \item $K_7$ – soziale Kontinua mit Stabilitätszyklen auf Gruppenebene;
    \item $K_8$ – zivilisatorische Kontinua mit großräumigen Koordinationsflüssen;
    \item $K_9$ – theoretische Kontinua (Struktur der Erklärungen);
    \item $K_{10}$ – metatheoretische Kontinua (Struktur struktureller Räume);
    \item $K_{11}$ – formale Ontologien und Kontinua auf logischer Ebene;
    \item $K_{12}$ – metaontologische Kontinua, die den Möglichkeitsraum regeln
                     aller unteren Ebenen.
\end{itemize}

Höhere Stufen ersetzen niedrigere nicht; sie betten sie ein und erweitern sie:
\[
K_x \subseteq M_x,
\qquad
M_x \subseteq \Omega(K_{x+1}).
\]


% ================================================================
\subsection{Grenzen, Zusammenbruch und Tod von Continua}
% ================================================================

Jedes $K_x$ hat eine Grenze $\partial\Omega(K_x)$, wobei die
Kontinuum verliert seine Zulässigkeit.
Ein Kontinuum stirbt, wenn:
\[
\Omega(K_x) = \varnothing
\quad\Rightarrow\quad
k(K_x) = 0,
\]
was entspricht:
\begin{itemize}
    \item Verletzung der Schwellenwerte $\Theta(K_x)$,
    \item Verlust stabiler Zyklen $C(K_x)$,
    \item destruktive Strömungen, die kritische Werte erreichen,
    \item Zusammenbruch der Achsen oder Inkompatibilität mit $M_x$.
\end{itemize}

Der Tod unterscheidet sich von der Evolution: Er entfernt das Kontinuum aus dem
den Raum der zulässigen Strukturen und kann nicht rückgängig gemacht werden.


% ================================================================
\subsection{Rolle der \texorpdfstring{$k$}{k}-Ebenenstruktur im Kernmodell}
% ================================================================

Die $K$-Ebenen bilden das vertikale Rückgrat der Ontologie von
Fortsetzung.
Sie:
\begin{itemize}
    \item organisieren den Übergang vom physischen zum kognitiven zum sozialen Zustand
          metaontologische Strukturen;
    \item definiert die Rekursion $K_x \to M_x \to K_{x+1}$;
    \item stellt die für scoped notwendige Dimensionsleiter bereit
          Anwendbarkeit;
    \item schränkt die Entwicklung von Kontinuen durch Zulässigkeit ein
          Bedingungen;
    \item vereint alle Bereiche der Realität unter einem mathematischen Rahmen.
\end{itemize}

Die Struktur der $K$-Ebenen macht die Ontologie sowohl hierarchisch als auch
rekursiv: Jede Ebene hängt vom darüber liegenden Metaraum ab und schränkt diese ein
die Ebene darunter.
